\section{第~\ref{sec:dishbrain}~章　DishBrain}

试验台的结构和游戏结果取自一项实验工作及其后续研究；噪声公式和向好设置漂移的公式在本章推导，并用本书的模型做了检验，而培养皿中的学习机制尚未确定。

{\footnotesize
\setlength{\tabcolsep}{3pt}\renewcommand{\arraystretch}{1.15}
\begin{longtable}{>{\raggedright}p{0.5cm}>{\raggedright}p{4.0cm}>{\raggedright}p{5.6cm}>{\raggedright}p{2.3cm}>{\raggedright\arraybackslash}p{1.7cm}}
\toprule
序号 & 论断 & 实验及测得的结果 & 来源 & 状态 \\
\midrule\endfirsthead
\toprule
序号 & 论断 & 实验及测得的结果 & 来源 & 状态 \\
\midrule\endhead
1 & 试验台：每块芯片约 $800\,000$ 个小鼠皮层细胞或约 $10^6$ 个人类细胞；$8$~mm$^2$ 上 $26\,000$ 个电极，最多记录 $1024$ 个，经 $8$ 个电极刺激 & 该工作的方法：MaxOne芯片，$8$~mm$^2$ 上 $26\,000$ 个铂电极，记录 $1024$ 个通道，每秒 $20\,000$ 个采样；技术上最多可刺激 $32$ 个电极，能独立刺激的是 $8$ 个；小鼠培养物约从第 $10$ 天起使用 & \cite{Kagan2022} & 已测量 \\
2 & 节拍 $10$~ms的回路：运动区的脉冲移动球拍（图~\ref{fig:dishloop}） & 脉冲按 $10$~ms窗口计数；从脉冲到响应刺激的延迟约 $5$~ms，由硬件缓冲决定 & \cite{Kagan2022} & 已测量 \\
3 & 输入：位置编码（$8$ 个电极中的哪一个）和 $4$--$40$~Hz的频率编码 & 主实验中，刺激频率从球在远端墙边时的 $4$~Hz线性增加到球拍处的 $40$~Hz；球的刺激脉冲幅度为 $75$~mV；脉冲为双相矩形脉冲 & \cite{Kagan2022} & 已测量 \\
4 & 输出 $\Delta p=c\,(n_\uparrow-n_\downarrow)$~\eqref{eq:dishreadout}，窗口计数噪声 $1/\sqrt{RT}$ & 该工作描述了两区差值规则（各区按活动加权），但没有给出确切公式。噪声估计是泊松计数的推论，在本章推导 & \cite{Kagan2022}；本章推导 & 理论 \\
5 & 老师：击中后给可预测刺激 $75$~mV、$100$~Hz、$0.1$~s；未击中后给不可预测刺激 $150$~mV、$5$~Hz、$4$~s，随机位置和时刻，随后 $4$~s停顿 & 方法：击中后 $75$~mV、$100$~Hz、$100$~ms，同时施加到全部 $8$ 个电极；未击中后 $150$~mV、$5$~Hz，在 $4$~s内于随机时刻施加到随机电极，随后 $4$~s无刺激。培养物中没有多巴胺 & \cite{Kagan2022} & 已测量 \\
6 & 网络的行为是为了让自己的输入变得可预测 & 自由能原理是一个理论设想，作者据此设计了方案；实验与它相符，但不能把它与更简单的机制区分开（第~7--8~行） & \cite{Friston2010,Kagan2022} & 假说 \\
7 & 若失败晃动网络而成功不晃动，停留在某设置的时间 $\rho\propto1/P_{\text{miss}}$~\eqref{eq:dishdensity}（命题~\ref{prop:dishscramble}） & 由对称扰动马尔可夫链中的流量平衡得出，在本章推导。没有在培养物上直接检验 & 本章推导 & 理论 \\
8 & “未击中即打乱”模型会学习：击中比例 $0.21\to0.38$；每个回合后都扰动时 $\approx0.12$，无扰动时 $0.19$（图~\ref{fig:dishmodel}） & 计算，脚本 \texttt{models/dishbrain\_model.py}，$40$ 个模型培养物，各 $2000$ 个回合：$0.21\to0.38$；$0.13\to0.11$；$0.19\to0.19$ & --- & 模型 \\
9 & 只有处在完整回路中的活培养物，连续击球回合才会变长；对照组不学习 & $399$ 次 $20$~min的会话：小鼠培养物 $9$ 个（$101$ 次会话），人类 $11$ 个（$138$），只有培养基的芯片 $6$ 块（$80$），静息培养物 $20$ 个（$42$），仿真 $3$ 个（$38$）。只有活培养物在后 $15$~min的平均回合长度大于前 $5$~min；非神经元细胞（HEK293T）和只有培养基的芯片没有效应 & \cite{Kagan2022} & 已测量 \\
10 & 一次会话内的显著增长只出现在完整反馈下；以安静代替刺激时，会话结束时的表现好于无反馈，但效应较小 & $3$ 天内 $486$ 次会话：“刺激”（$n=27$）、“安静”（$17$）、“无反馈”（$15$）。随时间的增长只在“刺激”条件下显著；会话结束时，“安静”条件优于“无反馈”，但效应较小 & \cite{Kagan2022} & 已测量 \\
11 & 效应不大；网络是否能长期学习是一个开放问题 & 一次会话内的改进是可靠的，但据作者本人所说，几天里跨会话的学习并未稳定出现 & \cite{Kagan2022} & 已测量 \\
12 & 游戏时网络趋近临界状态，越接近，游戏打得越好 & 对同一批培养物的活动雪崩分析：结构化输入使网络趋近临界状态，游戏水平与接近程度相关；但单有临界性、而没有关于行动后果的信息，不足以产生学习 & \cite{Habibollahi2023} & 已测量 \\
13 & 培养物的“感知能力”未被证明 & 三十位研究者的回应文章：闭环中的行为变化既不证明智能，也不证明感受；没有能表明这一点的实验 & \cite{Balci2023} & 假说 \\
14 & 建议的第四个对照：未击中后施加时长和能量相同的可预测刺激（第~\ref{sec:dishspec}~节） & 这样的实验尚未做过；这是本书的建议 & --- & 假说 \\
\bottomrule
\end{longtable}}
